\lysh{16144_SOLUTION}{1}
% TODO: TikZ σχήμα
\noindent \textbf{ΛΥΣΗ}

\medskip

\begin{enumerate}[\textalpha)]
    \item $\vec{AB} \cdot \vec{A\Delta} = |\vec{AB}| \cdot |\vec{A\Delta}| \cdot \sigma\upsilon\nu \widehat{(\vec{AB}, \vec{A\Delta})} = 4 \cdot 4 \cdot \sigma\upsilon\nu 60^\circ = 16 \cdot \frac{1}{2} = 8$.

    \item Τα διανύσματα $\vec{A\Delta}$ και $\vec{B\Gamma}$ είναι ομόρροπα οπότε σχηματίζουν γωνία $0^\circ$. Οπότε,
    \[ \vec{A\Delta} \cdot \vec{B\Gamma} = |\vec{A\Delta}| \cdot |\vec{B\Gamma}| \cdot \sigma\upsilon\nu \widehat{(\vec{A\Delta}, \vec{B\Gamma})} = 4 \cdot 4 \cdot \sigma\upsilon\nu 0^\circ = 16 \cdot 1 = 16 \]

    \item Οι διαγώνιες του ρόμβου διχοτομούνται, οπότε $\vec{AO} = \vec{O\Gamma}$, και τέμνονται κάθετα, οπότε τα διανύσματα $\vec{O\Delta}$ και $\vec{O\Gamma}$ σχηματίζουν γωνία $90^\circ$ και έχουν εσωτερικό γινόμενο μηδέν. Συνεπώς $\vec{O\Delta} \cdot \vec{AO} = \vec{O\Delta} \cdot \vec{O\Gamma} = 0$.

    \item Το τρίγωνο $AB\Delta$ είναι ισοσκελές αφού οι δύο πλευρές του είναι πλευρές του ρόμβου, με γωνία της κορυφής $60^\circ$. Άρα το τρίγωνο $AB\Delta$ είναι ισόπλευρο με πλευρά $4$. Το $O$ είναι το σημείο τομής των διαγωνίων του, οπότε $BO = O\Delta = 2$.
    \[ \vec{O\Delta} \cdot \vec{OB} = |\vec{O\Delta}| \cdot |\vec{OB}| \cdot \sigma\upsilon\nu \widehat{(\vec{O\Delta}, \vec{OB})} = 2 \cdot 2 \cdot \sigma\upsilon\nu 180^\circ = 4 (-1) = -4 \]

    \item Για τη γωνία $\widehat{(\vec{A\Delta}, \vec{\Gamma\Delta})}$ ισχύει : $\widehat{(\vec{A\Delta}, \vec{\Gamma\Delta})} = \widehat{(\vec{A\Delta}, \vec{BA})} = 180^\circ - 60^\circ = 120^\circ$. Οπότε,
    \[ \vec{A\Delta} \cdot \vec{\Gamma\Delta} = |\vec{A\Delta}| \cdot |\vec{\Gamma\Delta}| \cdot \sigma\upsilon\nu \widehat{(\vec{A\Delta}, \vec{\Gamma\Delta})} = 4 \cdot 4 \cdot \sigma\upsilon\nu 120^\circ = 16 \left( -\frac{1}{2} \right) = -8 \]
\end{enumerate}
\end{lysh}